\subsection{Normal endomorphisms}

DEFINITION 4. --- Let E be a hilbertian space and \(u \in \mathcal{L}(\mathrm{E})\) . We say that \(u\) is normal if it commutes with its adjoint \(u^{*}\) .

For example, every automorphism \(u\) of the hilbertian space \(E\) is normal since we have \(uu^{*} = u^{*}u = 1_{E}\) .

PROPOSITION 7. --- For \(u \in \mathcal{L}(\mathrm{E})\) to be normal, it is necessary and sufficient that \(\| u(x) \| = \| u^{*}(x) \|\) for all \(x \in \mathrm{E}\) .

We define a hermitian form \(\Phi\) on \(E\) by

\[
\Phi (x, y) = \langle u u ^ {*} (x) | y \rangle - \langle u ^ {*} u (x) | y \rangle .
\]

For u to be normal, it is necessary and sufficient that \(\Phi = 0\) . By the polarization formulas (V, p.~2), this is equivalent to \(\Phi(x, x) = 0\) for all \(x \in E\) . The proposition now follows since

\[
\Phi (x, x) = \left\| u ^ {*} (x) \right\| ^ {2} - \left\| u (x) \right\| ^ {2}.
\]

PROPOSITION 8. --- Suppose that \(u \in \mathcal{L}(\mathrm{E})\) is normal. Let \(\mathbf{N}\) be the kernel of \(u\) and \(\mathbf{M}\) the orthogonal of \(\mathbf{N}\) in \(\mathrm{E}\) ; let \(m\) and \(n\) be two positive integers such that \(m + n \geqslant 1\) . Then \(\mathbf{N}\) is the kernel of \(u^{m}(u^{*})^{n}\) and \(\mathbf{M}\) is both the initial and the final subspace of \(u^{m}(u^{*})^{n}\) . In particular, \(\mathbf{M}\) is both the initial and the final subspace of \(u\) and of \(u^{*}\) , and is stable under \(u\) and \(u^{*}\) .

Prop. 7 shows that u and \(u^{*}\) have the same kernel N. By prop. 4, (ii) of V, p.~41, the subspace M of E is stable under u and \(u^{*}\) since this is so for \(N = M^{\circ}\) , since \(M \cap N = \{0\}\) , the endomorphisms of M induced by u and \(u^{*}\) are injective. Let \(v = u^{m}(u^{*})^{n}\) ; the preceding argument shows that the restriction of v to M (resp. N) is injective (resp. null), hence N is the kernel of v. Consequently, \(M = N^{\circ}\) is the initial subspace of v. By prop. 4, (i) of V, p.~41, the final subspace of v is equal to the initial subspace of \(v^{*}\) . But \(v^{*} = u^{n}(u^{*})^{m}\) and so the initial subspace of \(v^{*}\) is equal to M by the preceding.

COROLLARY. --- Let \(\lambda \in \mathbf{K}\) . The following subspaces of E are equal :

\begin{enumerate}
\def\labelenumi{\alph{enumi})}
\item
  the eigen subspace of \(u\) relative to \(\lambda\) ;
\item
  the eigen subspace of \(u^{*}\) relative to \(\overline{\lambda}\) ;
\item
  the primary subspace of \(u\) relative to \(\lambda\) (in other words, by LIE, VII, § 1, No.~1, the set of all vectors \(x\) of \(E\) for which there exists an integer \(n \geqslant 0\) such that \((u - \lambda .1_{E})^{n}(x) = 0\) );
\item
  the primary subspace of \(u^{*}\) relative to \(\overline{\lambda}\) .
\end{enumerate}

It is clear that \(w = u - \lambda \cdot 1_{E}\) is a normal endomorphism of E, hence the endomorphisms w, \(w^{*} = u^{*} - \overline{\lambda} \cdot 1_{E}\) , \(w^{n}\) and \((w^{*})^{n}\) of E have the same kernel by prop. 8.

